\documentclass[10pt,a4paper]{amsart}
\usepackage[margin=19mm]{geometry}
\usepackage{amsmath,amssymb,mathtools}
\usepackage[scr=rsfs,cal=euler]{mathalfa}
\usepackage{xfrac}
\usepackage[unicode,hidelinks,bookmarks=false]{hyperref}
\newcommand{\ie}{i.e.\ }
\newcommand{\cf}{cf.\ }
\newcommand{\resp}{respectively\ }
\newcommand{\Subsection}{\subsection}
\providecommand{\nobreakdash}{\nobreak\hbox{-}}
\setlength{\emergencystretch}{2em}
\multlinegap0pt
\newcommand{\skpt}{\par\smallskip}
\usepackage{enumerate}
\newcommand\ZZ{\mathbb{Z}}
\newcommand\QQ{\mathbb{Q}}
\newcommand\RR{\mathbb{R}}
\newcommand\cC{{\mathcal C}}
\newcommand\kk{\Bbbk}

\newcommand\til{\mathord\sim}
\newcommand\app{\mathord\approx}

\DeclareMathOperator\Ap{Ap}
\DeclareMathOperator\supp{supp}
\DeclareMathOperator\Betti{Betti}
\DeclareMathOperator\Span{span}

\DeclarePairedDelimiter\abs{\lvert}{\rvert}

\newtheorem{thm}{Theorem}[section]
\newtheorem{prop}[thm]{Proposition}
\newtheorem{lemma}[thm]{Lemma}
\newtheorem{coro}[thm]{Corollary}
\theoremstyle{definition}\newtheorem{defi}[thm]{Definition}
\title[Kunz fan chamber incidence]{Numerical semigroups, polyhedra, and posets IV: original Theorem5.5}
\author{Cole Brower, Joseph McDonough and Christopher O'Neill}
\date{Source: Algebraic Combinatorics8(2025),597--617; DOI10.5802/alco.425}
\begin{document}\maketitle
\noindent\textit{Editorial scope.} This unit retains exact original Theorem5.5 and its full three-part proof with complete necessary projected-face/factorization core. Supporting original results are uncounted. Finite illustrated examples, enumeration recipes and later applications are omitted. Original variable and index print slips remain literal. All data used in this general statement are in the retained text, so no external illustrative asset is reproduced.
\setcounter{section}{1}
\section{Background}
\label{sec:background}

In this section, we recall some necessary background on semigroups, polyhedral geometry, Kunz nilsemigroups, and the Kunz cone. For a more complete overview, see~\cite{numericalappl}, \cite{ziegler}, \cite[Section~2]{minprescard} and \cite[Section~3]{kunzfaces3}, respectively.  

For a commutative semigroup $(N, +)$, an element $\infty\in N$ is \emph{nil} if $a + \infty = \infty$ for all $a \in N$. We call an element $a \in N$ \emph{nilpotent} if $na = \infty$ for some $n \in \ZZ_{\geq 1}$, and \emph{partly cancellative} if $a + b = a + c \ne \infty$ implies $b = c$ for all $b,c \in N$. We say $N$ is a \emph{nilsemigroup} if $N$ has an identity element and every non-identity element is nilpotent, and that $N$ is \emph{partly cancellative} if every non-nil element of $N$ is partly cancellative.  Note that any nilsemigroup is \emph{reduced}, meaning its only unit is the identity.  We call any element of $N$ that cannot be written as the sum of two non-identity elements an \emph{atom} of $N$. 

All semigroups $N$ in this paper are assumed to be commutative, partly cancellative, finitely generated, and reduced.  Under these assumptions, the atoms $n_0, \ldots, n_k \in N$ comprise the unique minimal generating set~\cite{fingenmon}; we denote this by $N = \langle  n_0, \ldots, n_k \rangle$.  A \emph{factorization} of $n \in N$ is an expression
\[
n = z_0n_0 + \cdots + z_kn_k  
\]
where each $z_i \in \ZZ_{\geq 0}$. The \emph{set of factorizations} of $n \in N$ is the set
\[
\mathsf{Z}_N(n) = \{z \in \ZZ_{\geq 0}^{k+1} \mid n = z_0n_0 + \cdots + z_kn_k\} \subset \ZZ_{\geq 0}^{k+1}.
\]
The \emph{factorization homomorphism} $\varphi_N : \ZZ_{\geq 0}^{k+1} \to N$ is the semigroup homomorphism 
\[
\varphi_N(z_0,\ldots, z_k) = z_0n_0 + \cdots z_kn_k.
\]
The \emph{kernel} of $\varphi_N$ is 
\[
\ker \varphi_N = \left\{ (a,b) \in \ZZ_{\geq 0}^{k+1} \times \ZZ_{\geq 0}^{k+1} \mid \varphi_N(a) = \varphi_N(b) \right\}  
\]
which induces a \emph{congruence relation} $\til$ on $\ZZ_{\geq 0}^{k+1}$, setting $a \sim b$ whenever $(a,b) \in \ker \varphi_N$ (recall that a \emph{congruence} is an equivalence relation such that $a \sim b$ implies $a+c \sim b+c$ for any $a,b,c \in \ZZ_{\geq 0}^{k+1}$).  We call any such pair $(a,b)$ or $a \sim b$ a \emph{trade} of $N$.  
A set of trades $\rho$ is said to \emph{generate} $\til$ if $\til$ is the smallest congruence on $\ZZ_{\geq 0}^{k+1}$ containing $\rho$.  A~\emph{presentation} of $N$ is a set $\rho$ of trades obtained by taking a generating set for $\til$ and omitting any $a \sim b$ where $\varphi_N(a)$ is nil.  A presentation $\rho$ of $\til$ is \emph{minimal} if no proper subset of $\rho$ is a presentation of $\til$.  It is known that any two minimal presentations of a finitely generated, partly cancellative semigroup have the same cardinality~\cite{kunzfaces3,fingenmon}.  

The \emph{support} of a factorization $ z \in \ZZ_{\geq 0}^{k+1}$ is the set 
$$\supp(z) = \{i \mid z_i > 0\}$$
of nonzero coordinates.  For $Z \subseteq \ZZ_{\geq 0}^{k+1}$, define 
$$\supp(Z) = \{i \mid z_i > 0 \text{ for some } z \in Z\},$$
and the \emph{factorization graph} $\nabla_Z$, whose vertices are elements of $Z$, and two factorizations $z, z' \in Z$ are connected by an edge if $\supp(z) \cap \supp(z') \ne \emptyset$. For $n \in N$, we write $\nabla_n$ for the factorization graph whose vertex set is $\mathsf Z_N(n)$.  For each $i \in \supp(Z)$,~let
\[
Z - e_i = \{z-e_i \mid z \in Z, \, i \in \supp(z)\}.
\]

Suppose $N$ is a finite, partially cancellative nilsemigroup.  An \emph{outer Betti element} of~$N$ is a subset $B \subseteq \mathsf{Z}_N(\infty)$ such that
\begin{enumerate}[(i)]
\item for each $i \in \supp(B)$, $B - e_i = \mathsf{Z}_N(p)$ for some $p \in N \setminus \{\infty\}$, and
\item the graph $\nabla_B$ is connected.
\end{enumerate}

A \emph{numerical semigroup} $S$ is an additive subsemigroup of $(\ZZ_{\geq 0}, +)$ that is cofinite and contains $0$.  Numerical semigroups have a unique minimal generating set, the size of which we call the \emph{embedding dimension}, and the smallest element of which we call the \emph{multiplicity}. Letting $m$ be an element of $S$, the Ap\'ery set of $S$ is the set 
\[
\Ap(S;m) = \{n \in S \mid n - m \notin S\}
\]
containing the smallest element of $S$ in each equivalence class modulo $m$.  
Let $\approx$ denote the congruence on $S$ setting $a \approx b$ whenever $a = b$ or $a,b \notin \Ap(S;m)$.  The~quotient semigroup $S/\app$ is a nilsemigroup with one non-nil element for each element of $\Ap(S;m)$.  The \emph{Kunz nilsemigroup} of $S$ is given by $N = \ZZ_m \cup \{\infty\}$ as sets, and is obtained from $S/\app$ by replacing each non-nil element by its equivalence class in $\ZZ_m$.  

\begin{thm}[{\cite{kunzfaces3}}]\label{t:minprescard}
If $\rho$ is a minimal presentation of a numerical semigroup $S$, then~$|\rho| = |\rho'| + \beta$, where $\rho'$ is any minimal presentation for the Kunz nilsemigroup~$N$ of~$S$, and $\beta$ is the number of outer Betti elements of $N$.  
\end{thm}

\setcounter{thm}{2}
A \emph{rational polyhedral cone} $C \subseteq \RR^d$ is the set of solutions to a finite set of linear inequalities with rational coefficients, i.e.
\[
C = \{x \in \RR^d \mid Ax \ge 0\}
\]
for some rational matrix $A$.  We say $C$ is \emph{strongly convex} (or \emph{pointed}) if it does not contain any positive dimension linear subspace of $\RR^d$, and the \emph{dimension} of $C$ is the vector space dimension $\dim C = \dim_\RR \Span_{\RR}C$. If none of the rows of $A$ can be omitted without changing $C$, we call $A$ the \emph{H-description} or \emph{facet description} of $C$, and refer to each inequality therein as a \emph{defining inequality} or \emph{facet inequality} of $P$.  If $\dim C = d$, then the facet description is unique up to the reordering and scaling of rows.  
Given a facet inequality $a_1x_1 + \cdots + a_dx_d \geq 0$ of $C$, the intersection of $C$ with the hyperplane $a_1x_1 + \cdots + a_dx_d = 0$ is called a \emph{facet} of $C$. Any intersection of facets is called a \emph{face} of~$C$; note that the condition of strong convexity is equivalent to $\{0\}$ being a face of~$C$. Given a face $F \subseteq C$, the \emph{relative interior} of $F$, denoted $F^{\circ}$, is the set of points in $F$ that do not lie in any proper faces of $F$.
A finite collection $\mathcal F$ of polyhedral cones is called a \emph{polyhedral fan} if 
\begin{enumerate}[(i)]
\item 
for any $C \in \mathcal{F}$, every face of $C$ is also in $\mathcal{F}$, and

\item 
the intersection of any $C,D \in \mathcal{F}$ is a face of both and lies in $\mathcal F$.

\end{enumerate}
The elements of a fan $\mathcal F$ are often called its \emph{faces}.  
A fan is \emph{pure} if its maximal faces (under containment) have the same dimension, and in this case, we refer to the maximal faces as \emph{chambers}.  

For $m \geq 2$, the strongly convex cone $\cC_m \subseteq \RR^{m-1}_{\geq 0}$ with facet inequalities
\[
x_i + x_j \geq x_{i+j}
\qquad \text{for} \qquad
i,j \in \ZZ_m \setminus \{0\}
\qquad \text{with} \qquad
i + j \ne 0
\]
is called the \emph{Kunz cone}.  
A point $z = (z_1, \ldots, z_{m-1}) \in \cC_{m} \cap \ZZ^{m-1}$ is an \emph{Ap\'ery point} if each $z_i \equiv i \bmod m$.  
The following is the culmination of results from~\cite{kunzfaces3,kunzfaces1,kunz}.  

\begin{thm}\label{t:kunznilsemigroupface}
The Ap\'ery points of $\cC_m$ are in bijective correspondence with numerical semigroups containing $m$; more specifically, this bijection is given by $z \mapsto S$ where $\Ap(S;m) = \{0, z_1,\ldots, z_{m-1}\}$.  
Fix a face $F \subseteq \cC_m$, the set
$$H = \{0\} \cup \{i : x_i = 0 \text{ for all } x \in F\} \subseteq \ZZ_m$$
is a subgroup of $\ZZ_m$ (called the \emph{Kunz subgroup} of $F$).  Define a binary operation $\oplus$ on~$N = (\ZZ_m/H) \cup \{\infty\}$ so that $\infty$ is nil, and for any nonzero $a,b \in \ZZ_m$,
\[
\overline a \oplus \overline b = \begin{cases}
a+b & \text{if } x_a + x_b = x_{a+b} \text{ for all } x \in F; \\
\infty & \text{otherwise.}
\end{cases}
\]
Under the above definition, $(N, \oplus)$ is a partly cancellative nilsemigroup (called the \emph{Kunz nilsemigroup} of $F$), and if the Ap\'ery point of a numerical semigroup $S$ lies in $F$, then $H = \{0\}$ and the Kunz nilsemigroup of $S$ equals the Kunz nilsemigroup of $F$.  
\end{thm}

In view of the above theorem, we say a numerical semigroup $S$ lies in the face $F \subseteq \cC_m$, and write $S \in F$, if the Ap\'ery point corresponding to $S$ lies in $F$.  



\section{Kunz fans}
\label{sec:kunzfan}

Throughout this section, fix $A = \{p_1, \ldots, p_k\} \subseteq \ZZ_m \setminus \{0\}$ with $\gcd(A, m) = 1$, and let $p:\RR^{m-1} \to \RR^k$ denote the linear map that projects each point in $\RR^{m-1}$ onto the coordinates whose indices lie in $A$, i.e., 
$$p(y) = (y_{p_1}, \ldots, y_{p_k}).$$

\begin{defi}\label{d:kunzfan}
Let $\mathcal F(m;A)$ denote the set of faces $F \subseteq \cC_m$ for which each atom of the Kunz nilsemigroup $N$ of $F$ has a representative in $A$.  
The \emph{Kunz fan} of $A$ is the~set
$$
\mathcal G(m;A) = \{p(F) : F \in \mathcal F(m;A)\}
$$
of cones in $\RR^k$ (we prove in Theorem~\ref{t:purefan} that $\mathcal G(m;A)$ is indeed a fan).  Note that under this definition, $|A|$ may exceed the number of atoms of $N$, but this allowance is necessary to ensure $\mathcal G(m;A)$ is a fan; see Example~\href{https://alco.centre-mersenne.org/articles/10.5802/alco.425/}{3.3} for a discussion of this subtlety.  
\end{defi}

One of the main results of this section is that $\mathcal G(m;A)$ is pure, and that the maximal faces of $\mathcal G(m;A)$ with respect to containment are precisely those whose corresponding nilsemigroup lies in the following family.  

\begin{defi}\label{d:staircaseposet}
A finite, partly cancellative nilsemigroup $N$ is a \emph{staircase} nilsemigroup if every non-nil element uniquely factors as a sum of atoms, and a numerical semigroup $S$ is called \emph{staircase} if its Kunz nilsemigroup is staircase.  As each outer Betti element of a staircase nilsemigroup $N$ consists of a single factorization of $\infty$, when there can be no confusion we refer to each such factorization as an outer Betti element of $N$.  
\end{defi}

A numerical semigroup $S$ is staircase if and only if every element of $\Ap(S;m)$ factors uniquely.  
In this case, $S$ is said to have \emph{Ap\'ery set of unique expression}; such semigroups have been studied in the literature~\cite{rosalesapery} and were central to the constructions in~\cite{minprescard}.  

\setcounter{thm}{3}
Consider the piecewise linear map $q:\RR_{\ge 0}^k \to \RR_{\ge 0}^{m-1}$ given by $q(x) = y$, where
$$y_i = \min\{c \cdot x \mid c \in \ZZ_{\ge 0}^k \text{ with } c_1p_1 + \cdots + c_kp_k \equiv i \bmod m\}.$$

\begin{prop}\label{p:circleoflightsinverse}
Fix a face $F \in \mathcal F(m; A)$ with Kunz nilsemigroup $N$ and Kunz subgroup $H$.  The map $p$ is injective on $F$, and if $x \in p(F)$, then $y = q(x)$ satisfies
$$y_i = z_1 x_1 + \cdots + z_k x_k$$
for any factorization $z \in \mathsf Z_N(i)$.  In particular, $q$ restricts to a linear map on $p(F)$, and $q(p(y)) = y$ for every $y \in F$.  
\end{prop}

\begin{proof}
For each nonzero $i \in \ZZ_m$, Theorem~\ref{t:kunznilsemigroupface} implies one of the following must hold:\ (i) $i \in H$ and $y_i = 0$; (ii) $i \in a + H$ for some $a \in A$ and $y_i = y_a$; or (iii) for any factorization $z \in \mathsf Z_N(i)$, we have 
$$y_i = z_1 y_{p_1} + \cdots + z_k y_{p_k}$$
for all $y \in F$.  As such, for any $y, y' \in F$, if $p(y) = p(y')$, then $y_a = y_a'$ for all $a \in A$, and thus $y = y'$.  This proves $p$ is injective on $F$.  

Next, suppose $p(y) = x$, and let $y' = q(x)$.  Fix $i$, and suppose that $c \in \ZZ_{\ge 0}^k$ satisfies $y_i' = c \cdot x$.  We claim $c \in \mathsf Z_N(i)$.  Indeed, if $i = 0$, then $y_i = 0$, and if $i \in a + H$ for some $a \in A$, then $y_i = y_a$.  For all other cases, fix $j$ with $c_j > 0$.  This means $i + H$ covers $(i - p_j) + H$ in $N$, and by minimality of $c \cdot x$, $(c - e_j) \cdot x = y_{i - p_j}'$.  By induction, we may assume $c - e_j \in \mathsf Z_N(j)$, and so we have $c = (c - e_j) + e_j \in \mathsf Z_N(i)$ since $y_{i - p_j}' + y_{p_j} = y_i'$.  This proves the claim.  By the preceding paragraph, we now have $y' = y$, and the remaining claims all immediately follow.  
\end{proof}

It is not hard to show that $q(\RR_{\ge 0}^k) \subseteq \cC_m$, although the injectivity in Proposition~\ref{p:circleoflightsinverse} is lost if one considers input outside of the faces in $\mathcal G(m;A)$.  

We next consider the cone
$$C(m;A) = \{x \in \RR_{\ge 0}^k : x_i \le c \cdot x  \text{ for all } c \in \ZZ_{\ge 0}^k \text{ with } c_1p_1 + \cdots + c_kp_k \equiv p_i \bmod m\}.$$
Despite the fact that $C(m;A)$ is defined using an infinite collection of inequalities, only finitely many are necessary.  Indeed, if $c \in \ZZ_{\ge 0}^k$ has some $c_j \ge m$, then 
$$
c \cdot x \ge (c - me_j) \cdot x
\qquad \text{for all} \qquad
x \in \RR_{\ge 0}^k,
$$
so in the definition of $C(m;A)$ one may restrict to $c$ with coordinates in $[0, m-1]$.  In~particular, this means $C(m;A)$ is a rational polyhedral cone.  


\setcounter{thm}{6}
\begin{lemma}\label{l:interioratomsets}
For each $x \in C(m;A)^\circ$, the image $y = q(x)$ lies in a face $F \subseteq \cC_m$ with trivial Kunz subgroup, and the Kunz nilsemigroup $N$ of $F$ has atom set $A$.  
\end{lemma}

\begin{proof}
Since $C(m;A) \subset \RR_{\ge 0}^k$, each $x_i$ is positive, so $x$ has all positive coordinates as well.  This ensures $F$ has trivial Kunz subgroup.  Moreover, the definition of $q$ ensures $N$ has atom set contained in $A$, but if some $p_i \in A$ were not an atom of $N$, then for any $z \in \mathsf Z_N(p_i)$, we would have 
\begin{align*}
x_i
\ge y_{p_i}
&= z_1 y_{p_1} + \cdots + z_{i-1} y_{p_{i-1}} + z_{i+1} y_{p_{i+1}} + \cdots + z_k y_{p_k} \\
&= z_1 x_1 + \cdots + z_{i-1} x_{i-1} + z_{i+1} x_{i+1} + \cdots + z_k x_k,
\end{align*}
violating the assumption that $x$ lies in the interior of $C(m;A)$.  
\end{proof}

\begin{thm}\label{t:purefan}
The cones in $\mathcal G(m;A)$ form a polyhedral fan that is pure, and the union of the cones in $\mathcal G(m;A)$ equals $C(m;A)$.  
\end{thm}

\begin{proof}
Cleary, if a face $F \subseteq \cC_m$ lies in $\mathcal F(m;A)$, then all faces of $F$ must as well, so every face of a cone in $\mathcal G(m;A)$ is a cone in $\mathcal G(m;A)$.  Moreover, for any \mbox{$F, F' \in \mathcal F(m;A)$}, Proposition~\ref{p:circleoflightsinverse} implies $p(F) \cap p(F') = p(F \cap F')$, so the intersection of any two cones in $\mathcal G(m;A)$ is a cone in $\mathcal G(m;A)$.  This verifies $\mathcal G(m;A)$ is a polyhedral fan.  

Next, by Lemma~\ref{l:interioratomsets}, for each $x \in C(m;A)^\circ$, $y = q(x)$ lies in a face $F \in \mathcal F(m;A)$.   
As such, $p(y) = x$ by Proposition~\ref{p:circleoflightsinverse}, which lies in $p(F) \in \mathcal G(m;A)$.  This means the union of the cones in $\mathcal G(m;A)$ contains $C(m;A)^\circ$, and since rational polyhedral cones are topologically closed, the union must equal $C(m;A)$.  

The final claim to prove is that every maximal cone in $\mathcal G(m;A)$ has dimension $k$.  Indeed, since $\mathcal G(m;A)$ contains only finitely many cones, and their union equals the $k$-dimensional cone $C(m;A)$, the union of the $k$-dimensional cones in $\mathcal G(m;A)$ must also equal $C(m;A)$.  This means every cone in $\mathcal G(m;A)$ is contained in a $k$-dimenisional cone in $\mathcal G(m;A)$.  
\end{proof}

\begin{coro}\label{c:chamberstaircase}
A face $F \in \mathcal F(m;A)$ is a chamber if and only if its Kunz nilsemigroup is staircase.   
\end{coro}

\begin{proof}
A maximal face of $\mathcal F(m;A)$ has dimension $k$ by Theorem~\ref{t:purefan}, and since the number of atoms of its Kunz nilsemigroup $N$ is also $k$, \cite[Theorem~4.3]{kunzfaces3} implies~$N$ has no inner trades and therefore must be staircase.  
\end{proof}


\section{Outer Betti elements and facets}
\label{sec:facetouterbetti}

Throughout this section, let $A = \{p_1,\ldots,p_k\} \subset \ZZ_m \setminus \{0\}$ with $\gcd(A,m) = 1$, and let $N$ be a partly cancellative nilsemigroup with atom set $A$.

The main results of this section are Theorem~\ref{t:outerbettihdesc}, which gives an $H$-description of each face $F \subseteq \cC_m$ in terms of its corresponding Kunz nilsemigroup, and Theorem~\href{https://alco.centre-mersenne.org/articles/10.5802/alco.425/}{4.9}, which characterizes the finite, partly cancellative nilsemigroups that are~Kunz.  

\begin{defi}\label{d:modularnilsemigroup}
A \emph{modular}\footnote{no relation to modular lattices} nilsemigroup is a finite, partly cancellative nilsemigroup~$N$ together with a bijection $f : \ZZ_m \to N \setminus \{\infty\}$ such that for any $a,b\in\ZZ_m$, either $f(a) + f(b) = f(a+b)$ or $f(a) + f(b) = \infty$.  For convenience, we write $a \in \ZZ_m$ in place of $f(a) \in N$, effectively viewing $N = \ZZ_m \cup \{\infty\}$ as sets.  If $N = \langle f(p_1), \ldots, f(p_k) \rangle$ and~$z \in \ZZ_{\ge 0}^k$, we write $\overline z = f(z_1p_1 + \cdots + z_kp_k)$, so if $z \in \mathsf Z_N(i)$ for $i \ne \infty$, then $\overline z = i$.  

A \emph{Betti equality} of a modular nilsemigroup $N$ is an equation of the form 
\[
c \cdot x = c_1x_1 + \cdots + c_kx_k = c_1'x_1+ \cdots c_k'x_k = c'\cdot x  
\]
where $c,c' \in \mathsf{Z}_N(p)$ for some $p \in N \setminus \{\infty\}$.  Let $H_N$ equal to subspace of $\RR^k$ satisfying all such equalities for a given $N$ (this coincides with the nullspace of the presentation matrix defined in \cite[Section~4]{kunzfaces3}).  
A \emph{Betti inequality} of $N$ is an inequality of the form 
\[
z \cdot x = z_1x_1 + \cdots + z_kx_k \geq a_1x_1 + \cdots + a_kx_k = a \cdot x  
\]
where $z \in \mathsf{Z}_N(\infty)$ is an outer Betti factorization and $a \in \mathsf{Z}_N(\overline{z})$.  Let $F_N \subseteq H_N$ be the rational polyhedral cone containing points $x \in H_N$ that satisfy all Betti inequalities.  
\end{defi}

\setcounter{thm}{4}
In what follows, let
$$Z = \{z \in \mathsf{Z}_N(\infty) \mid z - e_i \notin \mathsf{Z}_N(\infty) \text{ for each } i \in \supp(z)\}$$
denote the set of minimal elements of $\mathsf Z_N(\infty)$ under the component-wise partial order.  

\begin{lemma}\label{l:locateouterbetti}
Suppose $N$ is a modular nilsemigroup and fix $x \in F_N^\circ$ and $y \in \mathsf{Z}_N(\infty)$.  There exist $b, c \in \ZZ_{\geq 0}^k$ such that $b$ is a factorization of an outer Betti element of $N$,
\[
y \cdot x = (b+c) \cdot x, \quad \text{and} \quad \overline{y} = \overline{b+c}.
\]
\end{lemma}

\begin{proof}
Choose $z \in Z$ so that $y = z + \ell$ for some $\ell \in \ZZ_{\geq 0}^k$.  If $z$ lies in an outer Betti element of $N$, then choosing $b = z$ and $c = \ell$ completes the proof.  Otherwise, there exists $y' \in \mathsf Z_N(\infty) \setminus Z$ such that $y' \sim z$.  Choose $z' \in Z$ so that $y' = z' + \ell'$ for some $\ell' \in \ZZ_{\geq 0}^k \setminus \{0\}$.  If $z'$ is a factorization of an outer Betti element of $N$, then choosing $b = z'$ and $c = \ell' + \ell$ completes the proof, since $y' \sim z$ ensures
$$
(b + c) \cdot x = (y' + \ell) \cdot x = (z + \ell) \cdot x = y \cdot x
\qquad \text{and} \qquad
\overline{b + c} = \overline{y' + \ell} = \overline{z + \ell} = \overline{y}.
$$
Otherwise, we may again repeat this process to obtain $z'' \in Z$.  Notice that 
\begin{align*}
(z \cdot x) - (z' \cdot x)
&= (y' \cdot x) - (z' \cdot x)
= \ell' \cdot x
\ge \min \{x_1, \ldots, x_k\},
\end{align*}
so this process must eventually terminate in a suitable choice of $b$ and $c$.  
\end{proof}

\begin{prop}\label{p:innerfactsminimal}
Suppose $N$ is modular, fix $p \in \ZZ_m \setminus \{\infty\}$, and fix $x \in F_N$ that satisfies all Betti inequalities strictly.  If $z \in \mathsf Z_N(p)$ and $z' \in \ZZ_{\ge 0}^k$ with $\overline{z'} = p$, then
\[
z'\cdot x \geq z \cdot x,
\]
with equality if and only if $z' \in \mathsf{Z}_{N}(p)$.  
\end{prop}

\begin{proof}
If $z' \in \mathsf{Z}_N(p)$, then the definition of $F_N$ ensures $z \cdot x = z'\cdot x$, so it suffices to show that if $z' \in \mathsf{Z}_N(\infty)$ with $\overline{z'} = p$, we have $z' \cdot x > z \cdot x$.  By Lemma~\ref{l:locateouterbetti}, there is a factorization $b$ of an outer Betti element and $c \in \ZZ_{\geq 0}^k$ such that 
\[
z' \cdot x = (b + c) \cdot x
\qquad \text{and} \qquad
\overline{b+c} = p.
\]
Fixing $a \in \mathsf{Z}_N(\overline{b})$, the Betti inequality $z' \cdot x > (a+c) \cdot x$ must hold.  If $a + c \in \mathsf{Z}_N(p)$, then $(a + c) \cdot x = z \cdot x$ and we are done.  Otherwise, $(a+c) \in \mathsf{Z}_N(\infty)$, and we can again apply Lemma~\ref{l:locateouterbetti} to obtain a factorization $b'$ of an outer Betti element, $c' \in \ZZ_{\ge 0}^k$, and $a' \in \mathsf{Z}_N(\overline{b_1})$ such that 
\[
z' \cdot x > (a+c) \cdot x = (b' + c') \cdot x > (a' + c')\cdot x
\qquad \text{and} \qquad
\overline{a+c} = \overline{a' + c'} = p.
\]
As in the proof of Lemma~\ref{l:locateouterbetti}, repeating this process eventually terminates in a factorization $a'' + c'' \in \mathsf Z_N(p)$, at which point we obtain $z' \cdot x > (a'' + c'') \cdot x = z \cdot x$.  
\end{proof}

\setcounter{thm}{7}
\begin{thm}\label{t:outerbettihdesc}
If $F \subseteq \cC_m$ is a face with Kunz nilsemigroup $N$, then $p(F) = F_N$.
\end{thm}

\begin{proof}
Proposition~\ref{p:circleoflightsinverse} implies $p(F) \subseteq F_N$.  Conversely, fix $x \in F_N$.  Since $p(F) \subseteq F_N$ and both are rational polyhedral cones, it suffices to assume $x \in F_N^\circ$.  Fix nonzero $p, q \in \ZZ_m$ so that $p + q \ne 0$, and fix factorizations $z_1 \in \mathsf{Z}_N(p)$ and $z_2 \in \mathsf{Z}_N(q)$. We~must show for any $z \in \mathsf{Z}_N(p + q)$, we have $(z_1+z_2)\cdot x \geq z \cdot x$. If $z_1 + z_2 \in \mathsf{Z}_N(p + q)$, then $(z_1 + z_2) \cdot x = z \cdot x$.  Otherwise, $z_1 + z_2 \in Z_N(\infty)$, so $(z_1 + z_2)\cdot x \geq z \cdot x$ by Lemma~\ref{p:innerfactsminimal}.  Thus $x \in p(F)$, thereby completing the proof.
\end{proof}

\section{Walking the Kunz fan}
\label{sec:kunzwalk}

Throughout this section, let $A = \{p_1,\ldots,p_k\} \subseteq \ZZ_m \setminus \{0\}$ with $\gcd(A,m) = 1$, and let $N$ be a Kunz nilsemigroup with atom set $A$.  

\begin{defi}\label{d:irredundantinequality}
Fix a staircase Kunz nilsemigroup $N$.  We say an outer Betti element $z \in \mathsf Z_N(\infty)$ is \emph{irredundant} if its Betti inequality $z \cdot x \ge a \cdot x$ defines a facet of $F_N$, where $a \in \mathsf Z_N(\overline z)$.  
\end{defi}

We briefly identify some redundant outer Betti elements, for use in Section~\href{https://alco.centre-mersenne.org/articles/10.5802/alco.425/}{6}.  

\begin{prop}\label{p:reduandantbettis}
Let $N$ be a staircase Kunz nilsemigroup.  
\begin{enumerate}[(a)]
\item 
Any two outer Betti elements $b, b' \in \mathsf Z_N(\infty)$ with $\overline b = \overline b'$ have disjoint support.  

\item 
If $b \in \mathsf Z_N(\infty)$ is an outer Betti element and $\overline b = 0$, then $b$ is irredundant if and only if $\supp(b) = \{i\}$ and $a_1, \ldots, a_k$ are distinct modulo $\gcd(i,m)$.  

\item 
There is at most one outer Betti element of $N$ with full support, and such an outer Betti element is redundant if $|A| > 1$.  

\end{enumerate}
\end{prop}

\begin{proof}
Suppose $b_i > 0$ and $b_i' > 0$ for some $i$.  Since $b$ and $b'$ are outer Betti elements, we have $b - e_i, b' - e_i \in Z_N(\overline b - p_i)$, and since $N$ is staircase, $b - e_i = b' - e_i$ so $b = b'$.  

Next, if $b_i > 0$ and $b_j > 0$, then $b$ is redundant by the Farkas lemma, as its Betti inequality follows from $x_i \ge 0$ and $x_j \ge 0$.  As such, suppose $b_i$ is the only nonzero entry of $b$.  Writing $d = \gcd(i,m)$ for the order of $i \in \ZZ_m$, \cite[Corollary~3.7]{kunzfaces1} and Proposition~\ref{p:circleoflightsinverse} imply $C(m; A) \cap \{x \in \RR^d : x_i = 0\}$ projects faithfully onto $C(d; \overline A \setminus \{\overline 0\})$, where $\overline A \subseteq \ZZ_d$ is the set of residue classes of $a_1, \ldots, a_k$ in $\ZZ_d$.  This cone has dimension $k-1$ if and only if $a_1, \ldots, a_k$ are distinct modulo $d$, thereby proving the second claim.  

Lastly, the final claim immediately follows from the first and second.  
\end{proof}

\begin{prop}\label{p:fanboundary}
Suppose $N$ is a staircase Kunz nilsemigroup and $b \in \mathsf Z_N(\infty)$ is an irredundant outer Betti element.  The~facet of $F_N$ supported by the Betti inequality for~$b$ lies on the boundary of $C(m;A)$ if and only if either (i)~$\overline b \in A$, or (ii)~$\overline b = 0$.  
\end{prop}

\begin{proof}
By definition, a facet inequality of $C(m;A)$ is either of the form $x_i \ge c \cdot x$ with 
$$c_1p_1 + \cdots + c_kp_k \equiv p_i \bmod m,$$
or of the form $x_i \ge 0$.  As such, an outer Betti element whose Betti inequality is one of these two forms must fall into the claimed case~(i) or~(ii), respectively.  
\end{proof}

\setcounter{thm}{4}
\begin{thm}\label{t:cutandpaste}
Suppose $\dim F_N = k - 1$ and $F_N$ is not on the boundary of $C(m;A)$.  
\begin{enumerate}[(a)]
\item 
There is a unique inner Betti element $p \in N$, and $\mathsf Z_N(p) = \{z, z'\}$.  

\item 
For any non-nil $q \in N$, there exist unique factorizations $w, w' \in \mathsf Z_N(q)$ such that $z \not\preceq w$ and $z' \not\preceq w'$, where $\preceq$ denotes the component-wise partial order on $\ZZ_{\ge 0}^k$.  Moreover, $w \ne w'$ if and only if $p \preceq q$ in $N$.  

\item 
For a staircase Kunz nilsemigroup $N'$ with $F_N \subseteq F_{N'}$, either (i) $z \not\preceq w$ for all $q \in N'$ and $w \in \mathsf Z_{N'}(q)$, or (ii) $z' \not\preceq w$ for all $q \in N'$ and $w \in \mathsf Z_{N'}(q)$.  

\end{enumerate}
\end{thm}

\begin{proof}
Fix $c \in \ZZ^k$ with $\gcd(c) = 1$ and $H_N = \{x \in \RR^k : c \cdot x = 0\}$.  Write $c = c^+ - c^-$ where $c^+, c^- \in \ZZ_{\ge 0}^k$ have disjoint support, and let $d \in \ZZ_{\ge 1}$ be minimal with $\overline{d c^+} = \overline{d c^-}$.  
For any Betti element $p \in N$ and factorizations $z, z' \in \mathsf Z_N(p)$, we have $z - z' \in \ZZ c$.  As~such, $|\mathsf Z_N(p)| = 2$, as only one factorization of $p$ can avoid the positive (or~negative) support of $c$.  In particular, $\mathsf Z_N(p) = \{z, z'\}$ such that $z \in d \ZZ_{\ge 1} c^+$ and $z' \in d \ZZ_{\ge 1} c^-$.  To prove~(a), it suffices to show $d \overline{c^+}$ is a Betti element of $N$.  However, this follows from Proposition~\ref{p:circleoflightsinverse} and the fact that $d \overline{c^+} = d \overline{c^-} \ne 0$ by Propositions~\ref{p:reduandantbettis}(c) and~\ref{p:fanboundary} since $F_N$ is not on the boundary of $C(m; A)$.  

Having proven part~(a), let $p \in N$ denote the unique inner Betti element and write $\mathsf Z_N(p) = \{z,z'\}$.  Part~(b) then follows from part~(a), since any two factorizations of a given non-nil element $q \in N$ differ by an integer multiple of $z - z'$, and Proposition~\ref{p:circleoflightsinverse} ensures that if a factorization is preceded by $z$ or $z'$, then performing the trade $z \sim z'$ or $z' \sim z$, respectively, results in another factorization of $q$.  Lastly, for any $x \in F_{N'}^\circ$, the factorizations of any non-nil $q \in N$ are totally ordered by their dot product with~$x$ (in fact, they form an arithmetic sequence with step size $(z - z') \cdot x$), so since $N'$ is staircase it must be as prescribed in part~(c).  
\end{proof}

\begin{thebibliography}{99}
\bibitem[1]{numericalappl} A.Assi,M.D'Anna and P.A.García-Sánchez, \emph{Numerical semigroups and applications}, second ed., RSME Springer Series3, Springer,Cham,[2020] \copyright2020.
\bibitem[32]{ziegler} G.M.Ziegler, \emph{Lectures on polytopes}, Graduate Texts in Mathematics152, Springer-Verlag,New York,1995.
\bibitem[9]{minprescard} C.Elmacioglu,K.Hilmer,C.O'Neill,M.Okandan and H.Park-Kaufmann, ``On the cardinality of minimal presentations of numerical semigroups'', \emph{Algebr. Comb.}\textbf{7}(2024),no.3,753--771.
\bibitem[13]{kunzfaces3} T.Gomes,C.O'Neill and E.Torres Davila, ``Numerical semigroups,polyhedra,and posets III: Minimal presentations and face dimension'', \emph{Electron. J. Combin.}\textbf{30}(2023),no.2,article2.57(25pages).
\bibitem[25]{fingenmon} J.C.Rosales and P.A.García-Sánchez, \emph{Finitely generated commutative monoids}, Nova Science Publishers,Commack,NY,1999.
\bibitem[18]{kunzfaces1} N.Kaplan and C.O'Neill, ``Numerical semigroups,polyhedra,and posets I: The group cone'', \emph{Comb. Theory}\textbf{1}(2021),article19(23pages).
\bibitem[19]{kunz} E.Kunz, \emph{Über die Klassifikation numerischer Halbgruppen}, Regensburger Mathematische Schriften11, Universität Regensburg,Fachbereich Mathematik,Regensburg,1987.
\bibitem[24]{rosalesapery} J.C.Rosales, ``Numerical semigroups with Apéry sets of unique expression'', \emph{J. Algebra}\textbf{226}(2000),no.1,479--487.
\end{thebibliography}
\end{document}
